\documentclass[11pt]{article}
\usepackage[a4paper,margin=1in]{geometry}
\usepackage{fontspec}
\usepackage[boldfont=false]{xeCJK}
\setCJKmainfont{FandolSong-Regular.otf}
\usepackage{amsmath}
\usepackage{amssymb}
\usepackage{array}
\usepackage{booktabs}
\usepackage{graphicx}
\usepackage{adjustbox}
\usepackage{enumitem}
\setlist[itemize]{topsep=2pt,partopsep=0pt,parsep=0pt,itemsep=2pt,leftmargin=1.4em}
\setlist[enumerate]{topsep=2pt,partopsep=0pt,parsep=0pt,itemsep=2pt,leftmargin=1.6em}
\usepackage{parskip}
\usepackage{fvextra}
\fvset{breaklines=true,breakanywhere=true,fontsize=\small}
\setlength{\parindent}{0pt}

\begin{document}

\begin{center}
  {\LARGE\bfseries Data Representation}\\[0.35em]
  {\large A Level Computer Science}
\end{center}
\vspace{0.6em}

\section*{User-defined data types}

The built-in types (\texttt{INTEGER}, \texttt{REAL}, \texttt{STRING}, \texttt{CHAR}, \texttt{BOOLEAN}) cover the simplest cases. For richer problems you can define a \textbf{user-defined type} 用户定义类型, making the code clearer and the compiler stricter.

\subsection*{Why they are needed}

A built-in \texttt{STRING} lets you store nonsense in a field that should hold one of a few legal values; a user-defined type can restrict it. Real entities are usually a \textbf{collection} of values of different types. And \texttt{DECLARE Taxi : Vehicle} is clearer (self-documenting) than \texttt{DECLARE Taxi : STRING}.

\subsection*{Non-composite types}

\subsubsection*{Enumerated type}

An \textbf{enumerated type} 枚举类型 has values that are a \textbf{fixed list of named constants}:

\begin{Verbatim}
TYPE Vehicle = (M100, M230, T101, T102, T120, T150)
DECLARE MyTaxi : Vehicle
MyTaxi ← T102
\end{Verbatim}

The names are values of the new type (stored internally as small integers); you cannot assign anything outside the list. Uses: days of the week, colours, status codes.

\subsubsection*{Pointer type}

A \textbf{pointer} 指针 holds the \textbf{memory address of another variable} (or \texttt{NULL} for "no target"). Pointers build dynamic structures (linked lists, trees) and pass references without copying.

\begin{Verbatim}
TYPE PNode = ^TNode    // pointer to a TNode
DECLARE p : PNode
p ← NEW TNode
p^.Value ← 42          // dereference to reach the fields
\end{Verbatim}

To \textbf{dereference} 解引用 (\texttt{p\textasciicircum{}}) means to reach the variable it points to.

\subsection*{Composite types}

A \textbf{composite type} 复合类型 groups several values under one name.

\begin{itemize}
\item \textbf{record} 记录 (Topic 10) — fields of different types in a \texttt{TYPE ... ENDTYPE} block.

\item \textbf{set} 集合 — an unordered collection of unique values, with operations add, remove, membership test, union, intersection:

\end{itemize}

\begin{Verbatim}
DECLARE Available : SET OF Colour
Available ← {Red, Blue}
IF Green IN Available THEN ...
\end{Verbatim}

\begin{itemize}
\item \textbf{class} 类 / \textbf{object} 对象 — the OOP composite type, combining data fields (\textbf{attributes} 属性) with operations on them (\textbf{methods} 方法). An object is an instance of a class:

\end{itemize}

\begin{Verbatim}
CLASS Taxi
    PRIVATE Capacity : INTEGER
    PUBLIC FUNCTION GetCapacity() RETURNS INTEGER
        RETURN Capacity
    ENDFUNCTION
ENDCLASS
\end{Verbatim}

\subsection*{Choosing a type}

Use \textbf{enumerated} for a value from a fixed list, \textbf{pointer} for indirection, \textbf{record} for a group of fields, \textbf{set} for an unordered unique collection, and \textbf{class} when you need state \textbf{and} behaviour together.

\section*{File organisation and access}

\textbf{File organisation} 文件组织 is how the data is laid out; the access method is how the program reaches a record.

\begin{itemize}
\item \textbf{serial file} 串行文件 — records in the \textbf{order added}, no sorting. Access is sequential only; appending is fast; searching is slow. Used for logs and audit trails.

\item \textbf{sequential file} 顺序文件 — records \textbf{sorted by a key}. Searching is faster (you can stop early or binary-search); inserting is slow (records must shift). Used for master files updated in batch.

\item \textbf{random (direct-access) file} 随机文件 — records at positions computed from the key (often by a hash). Direct access by key is very fast; reading in key order is harder. Used for large lookup tables and customer accounts.

\end{itemize}

\par\medskip
\noindent\begin{minipage}{\linewidth}
\centering
\adjustbox{max width=\linewidth, max totalheight=0.42\textheight}{%
  \includegraphics{/Users/joschka/Desktop/Programming/eXercise/output/eXam/handouts/a_level_computer_science/assets/13-p324-serial-file.pdf}}\\[0.3em]
\emph{Serial file: records are kept in the order they were added}
\end{minipage}\par\medskip

\par\medskip
\noindent\begin{minipage}{\linewidth}
\centering
\adjustbox{max width=\linewidth, max totalheight=0.42\textheight}{%
  \includegraphics{/Users/joschka/Desktop/Programming/eXercise/output/eXam/handouts/a_level_computer_science/assets/13-p325-sequential-file.pdf}}\\[0.3em]
\emph{Sequential file: records are sorted by a key field}
\end{minipage}\par\medskip

\par\medskip
\noindent\begin{minipage}{\linewidth}
\centering
\adjustbox{max width=\linewidth, max totalheight=0.42\textheight}{%
  \includegraphics{/Users/joschka/Desktop/Programming/eXercise/output/eXam/handouts/a_level_computer_science/assets/13-p325-random-file.pdf}}\\[0.3em]
\emph{Random file: records sit at positions computed from the key}
\end{minipage}\par\medskip

The two access methods are \textbf{sequential access} 顺序存取 (read from start to end) and \textbf{direct access} 直接存取 (jump straight to a known position). Match the structure to the dominant operation: single-key lookups favour random; in-order reports favour sequential.

\section*{Hashing}

A \textbf{hash function} 散列函数 takes a record key and produces an \textbf{address} where the record is stored. A good one is fast, \textbf{deterministic} 确定性, and spreads keys evenly.

Examples for $N$ slots: modulo hash \texttt{address ← key MOD N}; folding (split the key, add the pieces, MOD N); a string hash (sum the character codes, MOD N).

A \textbf{collision} 冲突 is when two keys hash to the same address. Resolutions:

\begin{itemize}
\item \textbf{linear probing} 线性探测 — try the next slot, wrapping around. Simple but clusters.

\item \textbf{chaining} 链接法 — each slot points to a \textbf{linked list} 链表 of records that hashed there.

\item rehashing — use a second hash function.

\end{itemize}

To search: hash the key, read that slot; if the keys match you are done, else follow the resolution strategy until a match or an empty slot. To insert: hash the key, write to that slot or the next free one. Keep the \textbf{load factor} 装填因子 (records ÷ slots) below about 70\% for near-O(1) lookups.

\section*{Floating-point numbers}

To store \textbf{real numbers} of very different sizes, computers use a \textbf{floating-point} 浮点 format — a binary form of scientific notation, with two fields:

\begin{itemize}
\item a \textbf{mantissa} 尾数 — the significant digits.

\item an \textbf{exponent} 指数 — the power of 2 to multiply by.

\end{itemize}

Both are stored as \textbf{two's complement} 补码 integers. The value is

\[
\text{number} = \text{mantissa} \times 2^{\text{exponent}}.
\]

The mantissa is a fixed-point fraction, with an implied binary point after the sign bit. So \texttt{0.1010000} means $1/2 + 1/8 = 0.625$, and with exponent \texttt{00000010} (= 2) the value is $0.625 \times 2^{2} = 2.5$.

\par\medskip
\noindent\begin{minipage}{\linewidth}
\centering
\adjustbox{max width=\linewidth, max totalheight=0.42\textheight}{%
  \includegraphics{/Users/joschka/Desktop/Programming/eXercise/output/eXam/handouts/a_level_computer_science/assets/13-p329-floating-point-format.pdf}}\\[0.3em]
\emph{The place values of an 8-bit mantissa and an 8-bit exponent}
\end{minipage}\par\medskip

\subsection*{Converting}

\begin{itemize}
\item \textbf{binary → denary}: read the mantissa (use two's-complement rules if negative) as a fraction, read the exponent as a signed integer, then multiply mantissa by $2^{\text{exponent}}$.

\item \textbf{denary → binary}: write the number as a binary fraction × a power of 2, then store the mantissa and exponent in the agreed formats.

\end{itemize}

\subsection*{Normalisation}

A number is \textbf{normalised} 规格化 when the first significant bit is \textbf{immediately after the binary point} (no wasted leading zeros). This \textbf{maximises precision}, because every mantissa bit carries information. To normalise, shift the mantissa left and decrease the exponent (or shift right and increase it) until the first significant bit is in place; the value is unchanged. For negative (two's-complement) mantissas, the sign bit (1) is followed immediately by a 0.

\subsection*{Approximation and rounding errors}

Many denary reals \textbf{cannot be stored exactly} in binary — e.g. $0.1_{10}$ is the repeating binary fraction $0.000110011\ldots_{2}$, which must be truncated. Consequences:

\begin{itemize}
\item \textbf{rounding errors} 舍入误差 build up over many operations (\texttt{0.1 + 0.2} is not exactly \texttt{0.3}).

\item \textbf{comparisons fail} — test \texttt{ABS(x - 0.3) \textless{} 1e-9} instead of \texttt{x = 0.3}.

\item subtracting two nearly-equal values loses precision.

\end{itemize}

For exact needs (currency), use \textbf{fixed-point} 定点 or \textbf{BCD} 二进码十进数 instead of floating-point.


\end{document}
